% GENERATED FILE. Edit canonical lessons, then run npm run generate:books.
% canonical-source-hash: fnv1a64:9263a9c1ed1998ac
% canonical-lessons: HI-C446-prakrti, HI-C446-paryavaran, HI-C446-paudha, HI-C446-khilna, HI-C446-kumhlana

\chapter{Nature}
\label{ch:hi-nature}
\hlchaptermodality{446}

\begin{chapteropening}
\textbf{By the end of this chapter:} \emph{I can talk about nature and the environment, and plants that bloom and plants that wilt.}

Complete the last lesson of chapter 446: I can talk about nature and the environment, and plants that bloom and plants that wilt.
\end{chapteropening}

\section[\texorpdfstring{\hi{प्रकृति} (prakṛti)}{प्रकृति (prakṛti)}]{\hi{प्रकृति} (prakṛti) — nature}

\label{lesson:HI-C446-prakrti}



\begin{quote}
Before the new one: say the Hindi for the high tide, then the Hindi for a shell.
\end{quote}

\subsection*{You'll want to know: \hi{प्रकृति}}
\textbf{\hi{प्रकृति}} — \emph{prakṛti} — \textquotedblleft{}nature\textquotedblright{}.

\begin{quote}
\hi{पहाड़ों} \hi{में} \hi{प्रकृति} \hi{बहुत} \hi{सुंदर} \hi{है।}
\end{quote}

\begin{quote}
\emph{pahāṛoṁ meṁ prakṛti bahut sundar hai.}
\end{quote}

\begin{quote}
Nature is very beautiful in the mountains.
\end{quote}

Feminine.

\begin{cousinweb}
\textbf{\hi{प्रकृति}} (\emph{prakṛti}) is a tatsama, Sanskrit \emph{prakṛti}, \textquotedblleft{}the original or natural form of a thing\textquotedblright{} (Monier-Williams); it also means a person's nature, their temperament.
\end{cousinweb}

\begin{grammarlens}[title={What you've built}]
One more word for getting around.
\end{grammarlens}

\subsection*{Your turn}
\begin{itemize}
\raggedright
  \item \emph{Say it:} \emph{prakṛti}
  \item \emph{Say it:} \emph{prakṛti}, once more
  \item \emph{Say it:} \emph{pahāṛoṁ meṁ prakṛti bahut sundar hai.} — Nature is very beautiful in the mountains.
  \item \emph{Recall:} say the Hindi for the high tide, then the Hindi for a shell, then say \emph{prakṛti} again
\end{itemize}

\begin{tcolorbox}[breakable,colback=teal!4,colframe=teal!35!black,title={Before you move on}]
What does \hi{प्रकृति} mean? (\textquotedblleft{}Nature\textquotedblright{}.) Say it once more.
\end{tcolorbox}
\section[\texorpdfstring{\hi{पर्यावरण} (paryāvaraṇ)}{पर्यावरण (paryāvaraṇ)}]{\hi{पर्यावरण} (paryāvaraṇ) — the environment}

\label{lesson:HI-C446-paryavaran}



\begin{quote}
Before the new one: say the Hindi for a shell, then the Hindi for nature.
\end{quote}

\subsection*{You'll want to know: \hi{पर्यावरण}}
\textbf{\hi{पर्यावरण}} — \emph{paryāvaraṇ} — \textquotedblleft{}the environment\textquotedblright{}.

\begin{quote}
\hi{पेड़} \hi{काटना} \hi{पर्यावरण} \hi{के} \hi{लिए} \hi{बुरा} \hi{है।}
\end{quote}

\begin{quote}
\emph{peṛ kāṭnā paryāvaraṇ ke liye burā hai.}
\end{quote}

\begin{quote}
Cutting down trees is bad for the environment.
\end{quote}

Masculine.

\begin{cousinweb}
\textbf{\hi{पर्यावरण}} (\emph{paryāvaraṇ}) is a modern coinage from Sanskrit \emph{pari}, \textquotedblleft{}around\textquotedblright{}, and \emph{āvaraṇa}, \textquotedblleft{}covering\textquotedblright{} (Monier-Williams): what covers us all round.
\end{cousinweb}

\begin{grammarlens}[title={What you've built}]
One more word for getting around.
\end{grammarlens}

\subsection*{Your turn}
\begin{itemize}
\raggedright
  \item \emph{Say it:} \emph{paryāvaraṇ}
  \item \emph{Say it:} \emph{paryāvaraṇ}, once more
  \item \emph{Say it:} \emph{peṛ kāṭnā paryāvaraṇ ke liye burā hai.} — Cutting down trees is bad for the environment.
  \item \emph{Recall:} say the Hindi for a shell, then the Hindi for nature, then say \emph{paryāvaraṇ} again
\end{itemize}

\begin{tcolorbox}[breakable,colback=teal!4,colframe=teal!35!black,title={Before you move on}]
What does \hi{पर्यावरण} mean? (\textquotedblleft{}The environment\textquotedblright{}.) Say it once more.
\end{tcolorbox}
\section[\texorpdfstring{\hi{पौधा} (paudhā)}{पौधा (paudhā)}]{\hi{पौधा} (paudhā) — a plant; a sapling}

\label{lesson:HI-C446-paudha}



\begin{quote}
Before the new one: say the Hindi for nature, then the Hindi for the environment.
\end{quote}

\subsection*{You'll want to know: \hi{पौधा}}
\textbf{\hi{पौधा}} — \emph{paudhā} — \textquotedblleft{}a plant; a sapling\textquotedblright{}.

\begin{quote}
\hi{मेरे} \hi{कमरे} \hi{में} \hi{एक} \hi{छोटा} \hi{पौधा} \hi{है।}
\end{quote}

\begin{quote}
\emph{mere kamre meṁ ek choṭā paudhā hai.}
\end{quote}

\begin{quote}
There is a small plant in my room.
\end{quote}

Masculine.

\begin{cousinweb}
\textbf{\hi{पौधा}} (\emph{paudhā}) is a tadbhava of Sanskrit \emph{pravṛddha}, \textquotedblleft{}grown up, great\textquotedblright{} (CDIAL 8807), which lists Hindi \emph{paudhā}, \textquotedblleft{}young plants for transplanting\textquotedblright{}; Hindi widened it to any plant.
\end{cousinweb}

\begin{grammarlens}[title={What you've built}]
One more word for getting around.
\end{grammarlens}

\subsection*{Your turn}
\begin{itemize}
\raggedright
  \item \emph{Say it:} \emph{paudhā}
  \item \emph{Say it:} \emph{paudhā}, once more
  \item \emph{Say it:} \emph{mere kamre meṁ ek choṭā paudhā hai.} — There is a small plant in my room.
  \item \emph{Recall:} say the Hindi for nature, then the Hindi for the environment, then say \emph{paudhā} again
\end{itemize}

\begin{tcolorbox}[breakable,colback=teal!4,colframe=teal!35!black,title={Before you move on}]
What does \hi{पौधा} mean? (\textquotedblleft{}A plant; a sapling\textquotedblright{}.) Say it once more.
\end{tcolorbox}
\section[\texorpdfstring{\hi{खिलना} (khilnā)}{खिलना (khilnā)}]{\hi{खिलना} (khilnā) — to bloom; to open (of a flower)}

\label{lesson:HI-C446-khilna}



\begin{quote}
Before the new one: say the Hindi for the environment, then the Hindi for a plant.
\end{quote}

\subsection*{You'll want to know: \hi{खिलना}}
\textbf{\hi{खिलना}} — \emph{khilnā} — \textquotedblleft{}to bloom; to open (of a flower)\textquotedblright{}.

\begin{quote}
\hi{वसंत} \hi{में} \hi{बगीचे} \hi{में} \hi{फूल} \hi{खिलते} \hi{हैं।}
\end{quote}

\begin{quote}
\emph{vasant meṁ bagīce meṁ phūl khilte haiṁ.}
\end{quote}

\begin{quote}
In spring the flowers bloom in the garden.
\end{quote}

Past (it …): \hi{वह} \hi{फूल} \hi{सुबह} \hi{खिला}, \emph{vah phūl subah khilā}.

\begin{cousinweb}
\textbf{\hi{खिलना}} (\emph{khilnā}) goes back to \emph{*khiḍ}, \emph{*khil}, \textquotedblleft{}be open, expand\textquotedblright{}, forms CDIAL reconstructs (CDIAL 3882), which lists Hindi \emph{khilnā}, \textquotedblleft{}to open, bloom\textquotedblright{}. A face \hi{खिलता} too, when it lights up.
\end{cousinweb}

\begin{grammarlens}[title={What you've built}]
One more word for getting around.
\end{grammarlens}

\subsection*{Your turn}
\begin{itemize}
\raggedright
  \item \emph{Say it:} \emph{khilnā}
  \item \emph{Say it:} \emph{khilnā}, once more
  \item \emph{Say it:} \emph{vasant meṁ bagīce meṁ phūl khilte haiṁ.} — In spring the flowers bloom in the garden.
  \item \emph{Recall:} say the Hindi for the environment, then the Hindi for a plant, then say \emph{khilnā} again
\end{itemize}

\begin{tcolorbox}[breakable,colback=teal!4,colframe=teal!35!black,title={Before you move on}]
What does \hi{खिलना} mean? (\textquotedblleft{}To bloom; to open (of a flower)\textquotedblright{}.) Say it once more.
\end{tcolorbox}
\section[\texorpdfstring{\hi{कुम्हलाना} (kumhlānā)}{कुम्हलाना (kumhlānā)}]{\hi{कुम्हलाना} (kumhlānā) — to wilt; to fade}

\label{lesson:HI-C446-kumhlana}



\begin{quote}
Before the new one: say the Hindi for a plant, then the Hindi for to bloom.
\end{quote}

\subsection*{You'll want to know: \hi{कुम्हलाना}}
\textbf{\hi{कुम्हलाना}} — \emph{kumhlānā} — \textquotedblleft{}to wilt; to fade\textquotedblright{}.

\begin{quote}
\hi{पानी} \hi{के} \hi{बिना} \hi{पौधे} \hi{कुम्हला} \hi{गए।}
\end{quote}

\begin{quote}
\emph{pānī ke binā paudhe kumhlā gae.}
\end{quote}

\begin{quote}
Without water the plants wilted.
\end{quote}

Past (it …): \hi{वह} \hi{फूल} \hi{धूप} \hi{में} \hi{कुम्हलाया}, \emph{vah phūl dhūp meṁ kumhlāyā}.

\begin{cousinweb}
\textbf{\hi{कुम्हलाना}} (\emph{kumhlānā}) goes back to \emph{*kōmh}, \textquotedblleft{}wither\textquotedblright{}, a form CDIAL reconstructs (CDIAL 3524), which lists Hindi \emph{kumhlānā}, \textquotedblleft{}to fade, wither\textquotedblright{}.
\end{cousinweb}

\begin{grammarlens}[title={What you've built}]
One more word for getting around.
\end{grammarlens}

\subsection*{Your turn}
\begin{itemize}
\raggedright
  \item \emph{Say it:} \emph{kumhlānā}
  \item \emph{Say it:} \emph{kumhlānā}, once more
  \item \emph{Say it:} \emph{pānī ke binā paudhe kumhlā gae.} — Without water the plants wilted.
  \item \emph{Recall:} say the Hindi for a plant, then the Hindi for to bloom, then say \emph{kumhlānā} again
\end{itemize}

\begin{tcolorbox}[breakable,colback=teal!4,colframe=teal!35!black,title={Before you move on}]
What does \hi{कुम्हलाना} mean? (\textquotedblleft{}To wilt; to fade\textquotedblright{}.) Say it once more.
\end{tcolorbox}
